\section{Discussion}
\label{sec:implications}

\paragraph{From a single score to a conditional one.}
The literature's implicit model has been
$\text{risk}(x)=f(\text{model},\text{data})$: some one signal --
similarity, alignment, magnitude -- should rank vulnerable candidates for
any unlearner. Our results replace it with
$\text{risk}(x \mid \text{channel})$. The interaction is the right
evidence for this claim precisely because it is immune to the two
confounds that plague probe leaderboards: a uniformly better probe and a
uniformly more predictable objective both cancel in the double
difference. The asymmetry we observe -- proximity probes retain partial
signal on output-channel objectives while gradient probes go essentially
blind on RMU -- also explains why single-score studies reached
inconsistent conclusions: their conclusions depended on which objectives
happened to be in the roster. That the conditioning variable is trivial
to evaluate -- a static read of the loss definition -- is a feature, not
a limitation: the scientific content lies not in the taxonomy but in the
demonstrated law that this five-second label governs which predictor
family carries signal, by how much, and with which failure mode
(gradient probes do not merely weaken on the representation channel;
they go blind). A simple declared predicate with a non-obvious,
falsifiable consequence is precisely what makes the router deployable
and the claim preregistrable.

\paragraph{Operational picture.}
A practitioner declares the checkpoint, candidate scope, parameter block,
loss, and probe settings; reads the parent objective's channel off its
loss definition; and runs the matched probe. The profile's high-risk tail
localizes where post-unlearning measurement should concentrate, and it
allocates the protection budget before any audit is opened. The router
adds no tuning surface: the channel table is a lookup, and both probes are
forward-only at candidate scale. The profile does not decide whether a
request should be served by unlearning, retraining, or refusal, and a high
score does not mandate preservation; those remain application decisions.

\paragraph{Selection is not confirmation.}
A profile-selected pool may guide optimization, but improvement on that
pool is not evidence of preservation -- it rewards the method for
protecting what it was handed. All confirmatory outcomes are therefore
sealed-audit quantities, with both mean and tail reported because
concentrated damage hides in means.

\paragraph{What the score means.}
Each profile is conditional on the model, block, loss, precision, and
radius, and should be recomputed after material model changes. Randomized
sensitivity is an update-sensitivity prior rather than an
unlearning-specific quantity; its unlearning reading is licensed by the
specificity controls, and its trustworthiness by the fidelity gate rather
than by assumption.

\paragraph{Prediction and intervention are separate tests.}
A predictor can rank damage yet fail to support safe repair, and a repair
procedure can help for reasons unrelated to its selector. The crossed
design keeps the two claims apart: Table~\ref{tab:channel-matrix}
establishes conditional prediction, and Table~\ref{tab:crossed}
establishes that acting on the matched prediction -- and only the matched
prediction -- moves the sealed tail.
